\section{Conclusion}
\label{sec:conclusion}

I presented an end-to-end pipeline that builds and IRT-calibrates a
difficulty-graded question bank for biopharma regulatory compliance, and
that then uses the bank as a measuring instrument.

Part~1 built the instrument. It generates hard items under a Red Herring
plus subtle-violation mandate, administers them to a synthetic examinee
population, and freezes a healthy bank of $1{,}284$ Rasch-calibrated
items spanning roughly $\pm 4$ logits. A parallel track built from FDA
warning letters gives items that are harder than the synthetic ones
($+0.94$ logits, two-sided $p = 0.0012$), although the difference is a
shift in the lower tail rather than a full separation.

Part~2 ran two experiments on that instrument. Treating the calibrated
ability $\theta$ as an optimization metric, and using a strict and a
lenient judge on identical response text, I found that system-prompt
strictness primarily unlocks citation \emph{format} rather than
compliance \emph{reasoning}, with only the largest model showing a
genuine reasoning gain. Architectural interventions, namely retrieval,
self-critique, and step decomposition, showed the same pattern. All three
raised $\theta$ under lenient scoring and lowered it under strict
scoring. So the effect is not specific to prompt editing. When a judge
enforces section-level citations, citation format dominates the
measurement scale.

Part~3 turned the question toward the instrument itself. An audit against
two independent evaluation frameworks showed that a citation-requiring
G-Eval rubric recovers my judge's verdicts at $\mathrm{AUC} = 0.838$,
while a groundedness metric performs at chance. That is the pattern the
design predicts, and it rules out the possibility that the verdicts are
arbitrary. It does not show that they are correct, because every metric
in the audit shares a base model with the judge. The expert-review
protocol described in Section~\ref{sec:expertreview} is designed to close
that gap, and its labels are being collected now.

The remaining work follows from the limitations. The analytic difficulty
SE should be replaced with \texttt{py-irt} posterior standard deviations,
the examinee population should be expanded across model families to
reduce correlation, and the Red Herring mandate should be audited rather
than merely enforced. Beyond measurement, a calibrated bank enables
adaptive item selection, meaning the choice of the item whose difficulty
best matches a system's current $\theta$, and it gives a principled way
to track whether an intervention moves regulatory reasoning on a stable
scale.

The broader point is methodological. A single accuracy number cannot
distinguish a system that reasons better from one that has learned the
grader's formatting preference. Separating item difficulty from examinee
ability makes that distinction measurable, and in this domain the
distinction turned out to matter more than the headline numbers
suggested.
